% Practice list A v3 --- topics of Listas 1-2, QUALIFIER LEVEL (style guide §8: cross-topic,
% parameter analysis, graphs, answer key up to 15 pages).
% \solucoesfalse -> blank list ; \solucoestrue -> answer key + rubric.
% Compile: pdflatex (2x); check: pdffonts (all Type 1). Every number: practice-L1-L2-check.py
\documentclass[11pt,a4paper]{article}
% Shared preamble of the practice lists (practice-L*.tex), adapted from lista-exam.tex.
% Each list sets \solucoesfalse / \solucoestrue and \listatema right after % Shared preamble of the practice lists (practice-L*.tex), adapted from lista-exam.tex.
% Each list sets \solucoesfalse / \solucoestrue and \listatema right after % Shared preamble of the practice lists (practice-L*.tex), adapted from lista-exam.tex.
% Each list sets \solucoesfalse / \solucoestrue and \listatema right after \input{practice-preamble}.
\usepackage[utf8]{inputenc}
\usepackage[T1]{fontenc}
\usepackage{lmodern}   % Type 1 vector fonts (replaces PK bitmaps)
\usepackage{cmap}      % ToUnicode CMap -> copyable text and accents
\usepackage[english]{babel}
\usepackage{amsmath,amssymb,amsfonts}
\usepackage[a4paper,margin=2.2cm]{geometry}
\usepackage{enumitem}
\usepackage{fancyhdr}
\usepackage{comment}
\usepackage{xcolor}
\usepackage{microtype}
\usepackage{booktabs}
\usepackage{tikz}
\usepackage{pgfplots}
\pgfplotsset{compat=1.16}

\newif\ifsolucoes
\newcommand{\disciplina}{Macroeconomics I --- MPE Insper}
\newcommand{\espacoresol}{4.5cm}   % height of the answer space in the blank list
% The sol environment must be decided after the list sets its toggle.
\AtBeginDocument{\ifsolucoes\includecomment{sol}\else\excludecomment{sol}\fi}

\pagestyle{fancy}\fancyhf{}
\lhead{\footnotesize\disciplina}
\rhead{\footnotesize Practice list\ifsolucoes{} --- Solutions and rubric\fi}
\cfoot{\footnotesize\thepage}
\renewcommand{\headrulewidth}{0.3pt}
\setlength{\headheight}{14pt}
\setlist{topsep=2pt,itemsep=2pt}
\newcounter{questao}

\newcommand{\bloco}[3]{\bigskip\noindent
  {\large\bfseries\color{blue!55!black}Block #1 --- #2}\par\noindent\hfill\textit{(#3)}\par
  \noindent\rule{\linewidth}{0.3pt}\smallskip}
\newcommand{\questao}[1]{\refstepcounter{questao}\medskip\noindent
  \textbf{Question \thequestao.}~#1\par}
\newenvironment{alternativas}{\begin{enumerate}[label=\Alph*.,leftmargin=2.2em,topsep=1pt]}{\end{enumerate}}
\newcommand{\gabarito}[1]{\ifsolucoes\par\noindent\textbf{Answer:}~#1\fi}
\newcommand{\resolespaco}{\ifsolucoes\else\par\textit{Solution:}\vspace{\espacoresol}\par\fi}
\newcommand{\passo}[2]{\par\smallskip\noindent\textbf{Step --- #1.}~#2}
\newcommand{\intuicao}[1]{\par\smallskip\noindent\textit{Economic intuition.}~#1}
\newcommand{\refbib}[1]{\par\smallskip\noindent{\footnotesize\color{black!60}\textbf{Ref:} #1}}
\newcommand{\rubrica}[1]{\par\smallskip\noindent\textbf{Grading rubric.}\begin{itemize}[leftmargin=1.4em,topsep=1pt]#1\end{itemize}}

\newcommand{\cabecalho}{%
\begin{center}
  {\Large\bfseries \disciplina}\\[2pt]
  {\large Practice list: \listatema}\ifsolucoes\\[2pt]\textit{Solutions and grading rubric}\fi
\end{center}
\vspace{2pt}
{\footnotesize\noindent\textbf{Instructions.} Multiple choice: mark the single correct option (no
penalty for wrong answers). \emph{True/False}: every 2 wrong items cancel 1 right one
(score $=\max\{0,\text{right}-\text{wrong}/2\}\times$ item value); blank items do not count.
Discursive questions: show every step. Numbers are illustrative and chosen to give clean answers.\par}}
.
\usepackage[utf8]{inputenc}
\usepackage[T1]{fontenc}
\usepackage{lmodern}   % Type 1 vector fonts (replaces PK bitmaps)
\usepackage{cmap}      % ToUnicode CMap -> copyable text and accents
\usepackage[english]{babel}
\usepackage{amsmath,amssymb,amsfonts}
\usepackage[a4paper,margin=2.2cm]{geometry}
\usepackage{enumitem}
\usepackage{fancyhdr}
\usepackage{comment}
\usepackage{xcolor}
\usepackage{microtype}
\usepackage{booktabs}
\usepackage{tikz}
\usepackage{pgfplots}
\pgfplotsset{compat=1.16}

\newif\ifsolucoes
\newcommand{\disciplina}{Macroeconomics I --- MPE Insper}
\newcommand{\espacoresol}{4.5cm}   % height of the answer space in the blank list
% The sol environment must be decided after the list sets its toggle.
\AtBeginDocument{\ifsolucoes\includecomment{sol}\else\excludecomment{sol}\fi}

\pagestyle{fancy}\fancyhf{}
\lhead{\footnotesize\disciplina}
\rhead{\footnotesize Practice list\ifsolucoes{} --- Solutions and rubric\fi}
\cfoot{\footnotesize\thepage}
\renewcommand{\headrulewidth}{0.3pt}
\setlength{\headheight}{14pt}
\setlist{topsep=2pt,itemsep=2pt}
\newcounter{questao}

\newcommand{\bloco}[3]{\bigskip\noindent
  {\large\bfseries\color{blue!55!black}Block #1 --- #2}\par\noindent\hfill\textit{(#3)}\par
  \noindent\rule{\linewidth}{0.3pt}\smallskip}
\newcommand{\questao}[1]{\refstepcounter{questao}\medskip\noindent
  \textbf{Question \thequestao.}~#1\par}
\newenvironment{alternativas}{\begin{enumerate}[label=\Alph*.,leftmargin=2.2em,topsep=1pt]}{\end{enumerate}}
\newcommand{\gabarito}[1]{\ifsolucoes\par\noindent\textbf{Answer:}~#1\fi}
\newcommand{\resolespaco}{\ifsolucoes\else\par\textit{Solution:}\vspace{\espacoresol}\par\fi}
\newcommand{\passo}[2]{\par\smallskip\noindent\textbf{Step --- #1.}~#2}
\newcommand{\intuicao}[1]{\par\smallskip\noindent\textit{Economic intuition.}~#1}
\newcommand{\refbib}[1]{\par\smallskip\noindent{\footnotesize\color{black!60}\textbf{Ref:} #1}}
\newcommand{\rubrica}[1]{\par\smallskip\noindent\textbf{Grading rubric.}\begin{itemize}[leftmargin=1.4em,topsep=1pt]#1\end{itemize}}

\newcommand{\cabecalho}{%
\begin{center}
  {\Large\bfseries \disciplina}\\[2pt]
  {\large Practice list: \listatema}\ifsolucoes\\[2pt]\textit{Solutions and grading rubric}\fi
\end{center}
\vspace{2pt}
{\footnotesize\noindent\textbf{Instructions.} Multiple choice: mark the single correct option (no
penalty for wrong answers). \emph{True/False}: every 2 wrong items cancel 1 right one
(score $=\max\{0,\text{right}-\text{wrong}/2\}\times$ item value); blank items do not count.
Discursive questions: show every step. Numbers are illustrative and chosen to give clean answers.\par}}
.
\usepackage[utf8]{inputenc}
\usepackage[T1]{fontenc}
\usepackage{lmodern}   % Type 1 vector fonts (replaces PK bitmaps)
\usepackage{cmap}      % ToUnicode CMap -> copyable text and accents
\usepackage[english]{babel}
\usepackage{amsmath,amssymb,amsfonts}
\usepackage[a4paper,margin=2.2cm]{geometry}
\usepackage{enumitem}
\usepackage{fancyhdr}
\usepackage{comment}
\usepackage{xcolor}
\usepackage{microtype}
\usepackage{booktabs}
\usepackage{tikz}
\usepackage{pgfplots}
\pgfplotsset{compat=1.16}

\newif\ifsolucoes
\newcommand{\disciplina}{Macroeconomics I --- MPE Insper}
\newcommand{\espacoresol}{4.5cm}   % height of the answer space in the blank list
% The sol environment must be decided after the list sets its toggle.
\AtBeginDocument{\ifsolucoes\includecomment{sol}\else\excludecomment{sol}\fi}

\pagestyle{fancy}\fancyhf{}
\lhead{\footnotesize\disciplina}
\rhead{\footnotesize Practice list\ifsolucoes{} --- Solutions and rubric\fi}
\cfoot{\footnotesize\thepage}
\renewcommand{\headrulewidth}{0.3pt}
\setlength{\headheight}{14pt}
\setlist{topsep=2pt,itemsep=2pt}
\newcounter{questao}

\newcommand{\bloco}[3]{\bigskip\noindent
  {\large\bfseries\color{blue!55!black}Block #1 --- #2}\par\noindent\hfill\textit{(#3)}\par
  \noindent\rule{\linewidth}{0.3pt}\smallskip}
\newcommand{\questao}[1]{\refstepcounter{questao}\medskip\noindent
  \textbf{Question \thequestao.}~#1\par}
\newenvironment{alternativas}{\begin{enumerate}[label=\Alph*.,leftmargin=2.2em,topsep=1pt]}{\end{enumerate}}
\newcommand{\gabarito}[1]{\ifsolucoes\par\noindent\textbf{Answer:}~#1\fi}
\newcommand{\resolespaco}{\ifsolucoes\else\par\textit{Solution:}\vspace{\espacoresol}\par\fi}
\newcommand{\passo}[2]{\par\smallskip\noindent\textbf{Step --- #1.}~#2}
\newcommand{\intuicao}[1]{\par\smallskip\noindent\textit{Economic intuition.}~#1}
\newcommand{\refbib}[1]{\par\smallskip\noindent{\footnotesize\color{black!60}\textbf{Ref:} #1}}
\newcommand{\rubrica}[1]{\par\smallskip\noindent\textbf{Grading rubric.}\begin{itemize}[leftmargin=1.4em,topsep=1pt]#1\end{itemize}}

\newcommand{\cabecalho}{%
\begin{center}
  {\Large\bfseries \disciplina}\\[2pt]
  {\large Practice list: \listatema}\ifsolucoes\\[2pt]\textit{Solutions and grading rubric}\fi
\end{center}
\vspace{2pt}
{\footnotesize\noindent\textbf{Instructions.} Multiple choice: mark the single correct option (no
penalty for wrong answers). \emph{True/False}: every 2 wrong items cancel 1 right one
(score $=\max\{0,\text{right}-\text{wrong}/2\}\times$ item value); blank items do not count.
Discursive questions: show every step. Numbers are illustrative and chosen to give clean answers.\par}}

\solucoesfalse                 % <<<< \solucoestrue for the ANSWER KEY
\newcommand{\listatema}{Measurement, Welfare and Growth --- qualifier level (Listas 1--2)}
\renewcommand{\espacoresol}{6.5cm}

\begin{document}
\cabecalho

% =============================================================================
\bloco{1}{Price indices, new goods, welfare and growth accounting}{Kurlat chs.~1--2, 5 \,\textbullet\, 25 points}

\questao{(10 points) Objective items.}

\textbf{a)} (3 points) A household has CES preferences with elasticity of substitution $\eta=3$ between two
goods and, at base prices, spends half of its budget on each. The price of good 1 rises 50\% and the price
of good 2 does not change. By how much does a Laspeyres index overstate the true change in the cost of
living?
\begin{alternativas}
\item $2.1\%$
\item $6.2\%$
\item $25.0\%$
\item $0\%$
\end{alternativas}
\gabarito{B}
\begin{sol}
\passo{exact index}{For homothetic CES the cost of a util is $P=\big(\sum\omega_ip_i^{1-\eta}\big)^{1/(1-\eta)}$, so the
exact index is the power mean of order $1-\eta$ of the price relatives with base expenditure weights
(Question 2(b)): $\big(\tfrac12 1.5^{-2}+\tfrac12\big)^{-1/2}=0.7222^{-1/2}=1.1767$.}
\passo{bias}{Laspeyres $=\tfrac12\cdot1.5+\tfrac12=1.25$; $1.25/1.1767-1=6.2\%$. Option A is the Cobb--Douglas answer
($\eta=1$, geometric mean $\sqrt{1.5}$): the bias grows with substitutability. Second-order check:
$\tfrac\eta2\operatorname{Var}(\ln R)=1.5\cdot0.25(\ln1.5)^2=0.062$.}
\refbib{Kurlat (2020), \S1.2, pp.~22--25.}
\end{sol}

\medskip
\textbf{b)} (3 points) \emph{Welfare $\times$ growth.} Over 30 years Brazil's mean consumption grows 1.5\% a year, life
expectancy rises from 66 to 76 years, the variance of log consumption falls from 1.2 to 1.0, and leisure does
not change. With the Kurlat--Jones--Klenow utility, $\sigma=1$, $\bar u=5$ and initial mean consumption
normalized to one, the average annual growth of consumption-equivalent welfare is:
\begin{alternativas}
\item $1.50\%$
\item $1.83\%$
\item $4.33\%$
\item $2.30\%$
\end{alternativas}
\gabarito{C}
\begin{sol}
\passo{indifference across dates}{$\lambda$ scales initial consumption to match final welfare:
$a_0[\bar u+\ln\lambda+\mathbb E\ln c_0]=a_1[\bar u+\mathbb E\ln c_1]$, with $\mathbb E\ln c=\ln\bar c-v/2$ (lognormal):
$\mathbb E\ln c_0=-0.6$, $\mathbb E\ln c_1=0.45-0.5=-0.05$.}
\passo{solve}{$\ln\lambda=\frac{76}{66}(4.95)-5+0.6=1.30$, i.e.\ $1.30/30=4.33\%$ a year: consumption $0.45$,
inequality $0.10$, longevity $(76/66-1)\cdot4.95=0.75$. B omits longevity; D adds $\ln(76/66)$ instead of
weighting the value of life.}
\refbib{Kurlat (2020), \S2.2, eqs.\ (2.2.1)--(2.2.2).}
\end{sol}

\medskip
\textbf{c)} (2 points, T or F) A Paasche price index can never exceed the corresponding Laspeyres index.
\gabarito{F}
\begin{sol}
\passo{counterexample}{Equal base shares, price relatives $(1.2,1)$, quantity relatives $(1.3,1)$ (demand shifts toward
the good whose price rises): Laspeyres $1.100$, Paasche $\frac{0.5\cdot1.2\cdot1.3+0.5}{0.5\cdot1.3+0.5}=1.113$. In general
$P/L=1+\operatorname{Cov}_w(R_p,R_q)/(\bar R_p\bar R_q)$ (Bortkiewicz): Paasche exceeds Laspeyres when price and quantity
relatives covary positively.}
\refbib{Kurlat (2020), \S1.2.}
\end{sol}

\medskip
\textbf{d)} (2 points, T or F) \emph{PPP $\times$ productivity growth.} Tradables obey the law of one price and the
nominal exchange rate is constant. If Brazil's tradable-sector productivity grows 2 points a year faster than
the US's while non-tradable productivity grows equally in both, Brazil's price level relative to the US rises
by $2\gamma$ points a year, $\gamma$ being the non-tradable share in consumption.
\gabarito{T}
\begin{sol}
\passo{Balassa--Samuelson in growth rates}{$P=p_T(A_T/A_N)^{\gamma}$ in each country (common wage, zero profits), and
$p_T^{BR}=Ep_T^{US}$; hence $d\ln\frac{P^{BR}}{EP^{US}}=\gamma[(g_T-g_N)^{BR}-(g_T-g_N)^{US}]=2\gamma$ points.}
\refbib{Kurlat (2020), \S1.2; Romer (2012), \S4.2.}
\end{sol}

% -----------------------------------------------------------------------------
\questao{(15 points) Discursive. \emph{Measurement $\times$ growth accounting.} Consumers have CES preferences over
the set $\Omega_t$ of varieties available at date $t$,
\[U=\Big(\sum_{i\in\Omega_t}(b_ix_i)^{\frac{\eta-1}{\eta}}\Big)^{\frac{\eta}{\eta-1}},\qquad\eta>1,\]
where $b_i>0$ is a quality weight. The statistical agency observes prices and expenditure shares and
builds its deflators from varieties present in both periods.}

\textbf{a)} (3 points) Find the minimum expenditure needed to obtain one unit of utility at prices $p$, and
each variety's expenditure share.
\resolespaco
\begin{sol}
\passo{problem}{$\min\sum_ip_ix_i$ s.t.\ $\sum_i(b_ix_i)^{\rho}=1$, $\rho\equiv\frac{\eta-1}{\eta}\in(0,1)$. FOC with multiplier $\mu$:
$p_i=\mu\rho b_i^{\rho}x_i^{\rho-1}$.}
\passo{demands}{$x_i=(\mu\rho)^{\eta}b_i^{\eta-1}p_i^{-\eta}$, using $\frac1{1-\rho}=\eta$ and $\frac{\rho}{1-\rho}=\eta-1$. The constraint
gives $(\mu\rho)^{\eta-1}\sum_i(p_i/b_i)^{1-\eta}=1$.}
\passo{cost of a util}{Spending $=(\mu\rho)^{\eta}\sum_ib_i^{\eta-1}p_i^{1-\eta}=(\mu\rho)^{\eta}(\mu\rho)^{1-\eta}=\mu\rho$, so
\[P=\Big(\sum_i(p_i/b_i)^{1-\eta}\Big)^{\frac1{1-\eta}},\qquad s_i=\frac{p_ix_i}{P}=\frac{(p_i/b_i)^{1-\eta}}{P^{1-\eta}}.\]
Homotheticity makes $P$ independent of the utility level: it is a true price index.}
\refbib{Kurlat (2020), \S1.2 (biases of price indices); formula sheet \texttt{rules/01}.}
\rubrica{
\item FOC and demands: 1.5 pts.
\item $P$ and the shares: 1.5 pts. A formula without derivation: max.\ 1 pt in the item.
}
\end{sol}

\textbf{b)} (3 points) Keep the set of varieties fixed between dates 0 and 1. Express the exact index $P_1/P_0$ in
terms of observable price relatives $R_i=p_{i1}/p_{i0}$ and shares, rank it against the Laspeyres and Paasche
indices, and approximate the Laspeyres bias to second order in the dispersion of $\ln R_i$.
\resolespaco
\begin{sol}
\passo{base weights}{$(P_1/P_0)^{1-\eta}=\sum_i(p_{i1}/b_i)^{1-\eta}/P_0^{1-\eta}=\sum_is_i^0R_i^{1-\eta}$: the power mean of order $1-\eta$
with base weights, $M_{1-\eta}(R;s^0)$.}
\passo{current weights}{$s_i^1=s_i^0R_i^{1-\eta}/\sum_js_j^0R_j^{1-\eta}$, so $\sum_is_i^1R_i^{\eta-1}=1/\sum_js_j^0R_j^{1-\eta}$ and
$P_1/P_0=M_{\eta-1}(R;s^1)$.}
\passo{ranking}{Laspeyres $=M_1(R;s^0)$, Paasche $=\big(\sum_is_i^1/R_i\big)^{-1}=M_{-1}(R;s^1)$. Power means increase with the
order; since $1-\eta<1$ and $\eta-1>-1$: Paasche $\le P_1/P_0\le$ Laspeyres, with equality iff all $R_i$ are equal.}
\passo{second order}{With $z_i=\ln R_i$, mean $\bar z$ and variance $V$ (base weights),
$\ln M_r=\frac1r\ln\sum s_ie^{rz_i}\approx\bar z+\frac r2V$. Hence $\ln L-\ln(P_1/P_0)\approx\frac12V-\frac{1-\eta}2V=\frac\eta2V$.}
\refbib{Kurlat (2020), \S1.2, pp.~22--25; Jones (2020), \S2.3.}
\rubrica{
\item Exact index with base and current weights: 1.5 pts.
\item Ranking by the power-mean inequality: 1 pt. Ranking asserted: 0 for this line.
\item Second-order bias: 0.5 pt.
}
\end{sol}

\textbf{c)} (3 points) Now some varieties are new at date 1 and none disappears. Find the exact index in terms
of the index over continuing varieties and observable shares. What happens as $\eta\to\infty$ and as $\eta\to1$?
\resolespaco
\begin{sol}
\passo{split the index}{Let $C$ be the continuing set and $\lambda_t\equiv\frac{\sum_{i\in C}(p_{it}/b_i)^{1-\eta}}{\sum_{i\in\Omega_t}(p_{it}/b_i)^{1-\eta}}$, which by (a) is
the expenditure share of continuing varieties at $t$ ($\lambda_0=1$). Then $P_t^{1-\eta}=(P_t^C)^{1-\eta}/\lambda_t$.}
\passo{formula}{Raising to $1/(1-\eta)$: $P_t=P_t^C\lambda_t^{1/(\eta-1)}$, so
\[\frac{P_1}{P_0}=\frac{P_1^C}{P_0^C}\Big(\frac{\lambda_1}{\lambda_0}\Big)^{\frac1{\eta-1}}.\]}
\passo{limits}{$\eta\to\infty$: the exponent $\to0$ and the correction disappears (perfect substitutes add nothing
except through price). $\eta\to1^+$: the exponent $\to\infty$ and, with $\lambda_1<1$, the correction $\to0$: new
varieties become infinitely valuable.}
\refbib{Kurlat (2020), \S1.2 (new-goods bias).}
\rubrica{
\item $\lambda_t$ as an observable share: 1 pt.
\item Formula derived: 1.5 pts.
\item Both limits: 0.5 pt.
}
\end{sol}

\textbf{d)} (3 points) Continuing-variety consumption inflation is 5\% and new varieties take 4\% of date-1
consumption spending; $\eta=4$. New varieties appear only among consumption goods (share $s_C=0.8$ of GDP);
investment goods (share $s_I=0.2$) have their own deflator bias $b_I$ from unmeasured quality change. In a
steady state with $\alpha=1/3$, how is the Solow residual of growth accounting distorted? Compute it for
$b_I=0$, $b_I=b_C$ and $b_I=3b_C$, where $b_C$ is the consumption-deflator bias.
\resolespaco
\begin{sol}
\passo{consumption bias}{True consumption inflation $1.05\cdot0.96^{1/3}-1=3.58\%$: $b_C=1.42$ points a year.}
\passo{measured growth}{Real GDP growth is nominal growth minus the deflator, so it is understated by
$b_Y=s_Cb_C+s_Ib_I$. Capital is built by perpetual inventory from deflated investment, so in a steady state its
measured growth is understated by $b_I$.}
\passo{residual}{$g_A^m=(g_Y-b_Y)-\alpha(g_K-b_I)-(1-\alpha)g_L$, hence
\[g_A^m-g_A=-(s_Cb_C+s_Ib_I)+\alpha b_I.\]
$b_I=0$: $-1.14$ points; $b_I=b_C$: $-(1-\alpha)b_C=-0.95$; $b_I=3b_C$: $-0.57$.}
\intuicao{Understated output lowers measured TFP growth; understated capital raises it, because the
residual absorbs whatever inputs miss. Quality change in equipment is the embodied technical progress of
Romer's problem 1.12, which growth accounting books as TFP.}
\refbib{Kurlat (2020), \S5.4, eq.\ (5.4.2), p.~87; Romer (2012), \S1.7 and problem 1.12.}
\rubrica{
\item $b_C$: 0.5 pt.
\item Biases of measured output and capital growth: 1 pt.
\item Residual formula and the three numbers: 1.5 pts. Ignoring the capital channel: max.\ 1.5 pts in the item.
}
\end{sol}

\textbf{e)} (3 points) \emph{Parameter analysis.} How do the inflation bias and the residual bias depend on $\eta$, the
new-variety share $\nu\equiv1-\lambda_1$ and $b_I$? Prove the signs, fill the inflation bias for $\eta\in\{2,4,8\}$ and
$\nu\in\{2\%,4\%,8\%\}$, find the $b_I$ at which the residual bias changes sign, and draw the inflation bias against
$\eta$ for the three values of $\nu$.
\resolespaco
\begin{sol}
\passo{signs}{Bias $B(\eta,\nu)=1.05\,[1-(1-\nu)^{1/(\eta-1)}]$. Since $\ln(1-\nu)<0$, $(1-\nu)^{1/(\eta-1)}$ rises with $\eta$ and falls with
$\nu$: $\partial B/\partial\eta<0$, $\partial B/\partial\nu>0$. The residual bias $-(s_Cb_C)+(\alpha-s_I)b_I$ rises with $b_I$ iff
$\alpha>s_I$, and is zero at $b_I^*=s_Cb_C/(\alpha-s_I)=8.5$ points: only implausibly fast equipment quality change
would make measured TFP growth too high.}
\passo{table}{Inflation bias in points (continuing inflation 5\%):
\begin{center}
\begin{tabular}{lccc}
\toprule
 & $\nu=2\%$ & $\nu=4\%$ & $\nu=8\%$ \\
\midrule
$\eta=2$ & 2.10 & 4.20 & 8.40 \\
$\eta=4$ & 0.70 & 1.42 & 2.88 \\
$\eta=8$ & 0.30 & 0.61 & 1.24 \\
\bottomrule
\end{tabular}
\end{center}}
\begin{center}
\begin{tikzpicture}
\begin{axis}[width=10cm,height=5.2cm,xmin=1.5,xmax=10,ymin=0,ymax=17,
  xlabel={$\eta$},ylabel={bias (points)},legend style={font=\footnotesize},legend pos=north east,
  tick label style={font=\footnotesize},label style={font=\small}]
\addplot[thick,blue,domain=1.5:10,samples=120]{105*(1-(0.98)^(1/(x-1)))};
\addplot[thick,red,domain=1.5:10,samples=120]{105*(1-(0.96)^(1/(x-1)))};
\addplot[thick,black,dashed,domain=1.5:10,samples=120]{105*(1-(0.92)^(1/(x-1)))};
\legend{$\nu=2\%$,$\nu=4\%$,$\nu=8\%$}
\end{axis}
\end{tikzpicture}
\end{center}
\passo{limits and reading}{$\eta\to1$: $B\to105$ points, because the true index falls toward zero when new varieties
become essential; $\eta\to\infty$: $B\to0$. The bias is roughly proportional to $\nu$ and to $1/(\eta-1)$: it is a
problem for goods with many new, differentiated varieties, not for commodities.}
\refbib{Kurlat (2020), \S1.2 and \S5.4.}
\rubrica{
\item Signs proved: 1 pt. Signs listed without the argument: max.\ 0.5 pt here.
\item Table: 0.5 pt.
\item $b_I^*$ and its reading: 0.5 pt.
\item Graph (axes, three decreasing convex curves ordered by $\nu$): 1 pt.
}
\end{sol}

% =============================================================================
\bloco{2}{Growth, capital flows and national income}{Kurlat chs.~4--5; Romer chs.~1, 4 \,\textbullet\, 25 points}

\questao{(10 points) Objective items.}

\textbf{a)} (3 points) \emph{Development accounting $\times$ capital flows.} Let $Y=AK^{\alpha}(hL)^{1-\alpha}$ with $\alpha=0.4$ and the
same $A$ in Brazil and the US. Brazil's output per worker is 25\% of the US's and its human capital per worker
60\%. The ratio of marginal products of capital (Brazil/US) implied by the closed-economy model is about:
\begin{alternativas}
\item $8.0$
\item $3.7$
\item $1.0$
\item $2.5$
\end{alternativas}
\gabarito{B}
\begin{sol}
\passo{MPK in terms of observables}{$MPK=\alpha y/k$ and $k=(y/(Ah^{1-\alpha}))^{1/\alpha}$, so
$MPK=\alpha A^{1/\alpha}(y/h)^{(\alpha-1)/\alpha}$.}
\passo{ratio}{$(0.25/0.6)^{-1.5}=3.7$. Without human capital, $0.25^{-1.5}=8.0$ (A). Capital should flow to Brazil;
that it does not suggests lower $A$ (Lucas paradox).}
\refbib{Kurlat (2020), \S5.3 (the capital accumulation hypothesis), pp.~80--86.}
\end{sol}

\medskip
\textbf{b)} (3 points) With $f(k)=[\alpha k^{-1}+1-\alpha]^{-1}$ (elasticity of substitution $1/2$), $\alpha=1/2$ and $n+g+\delta=0.08$, the
Golden Rule saving rate is:
\begin{alternativas}
\item $50\%$
\item $20\%$
\item $8\%$
\item $32\%$
\end{alternativas}
\gabarito{B}
\begin{sol}
\passo{Golden Rule}{$f(k)=\frac{k}{\alpha+(1-\alpha)k}$, $f'(k)=\frac{\alpha}{(\alpha+(1-\alpha)k)^2}=0.08\Rightarrow\alpha+(1-\alpha)k=2.5$, $k=4$, $f=1.6$.}
\passo{saving rate}{$s=\frac{(n+g+\delta)k}{f(k)}=\frac{0.32}{1.6}=20\%$, which equals the capital share \emph{at} $k_{gr}$, not $\alpha$ (A).}
\refbib{Kurlat (2020), \S4.3; Romer (2012), \S1.8 (CES).}
\end{sol}

\medskip
\textbf{c)} (2 points, T or F) With CES production and an elasticity of substitution different from one, a
balanced growth path with a constant capital--output ratio exists only if technical progress is
labor-augmenting.
\gabarito{T}
\begin{sol}
\passo{Uzawa}{On a balanced path $K$ and $Y$ grow at the same rate. With $Y=F(BK,AL)$ CES,
$Y/K=[\alpha B^{\rho}+(1-\alpha)(AL/K)^{\rho}]^{1/\rho}$ is constant along an exact balanced path only if $B$ is constant
and $AL/K$ is constant, i.e.\ technology augments labor. Cobb--Douglas is the exception, where all forms are
equivalent.}
\refbib{Romer (2012), \S1.2 and problem 1.12.}
\end{sol}

\medskip
\textbf{d)} (2 points, T or F) \emph{Growth $\times$ open economy.} In a small open Solow economy with perfect capital
mobility, a permanent rise in the saving rate permanently raises GDP per worker.
\gabarito{F}
\begin{sol}
\passo{why}{GDP per effective worker is pinned by $f'(k)=r^*+\delta$; saving changes who owns the capital (net foreign
assets and national income), not domestic production (Question 4).}
\refbib{Romer (2012), problem 4.4.}
\end{sol}

% -----------------------------------------------------------------------------
\questao{(15 points) Discursive. \emph{Growth $\times$ national accounts $\times$ capital flows.} Brazil as a small open
economy. Output is $Y=K^{\alpha}(AhL)^{1-\alpha}$; capital moves freely and earns the world rate $r^*$ net of depreciation $\delta$;
$A$ grows at $g$, $L$ at $n$, $h$ is constant. Residents own wealth $W$ (capital at home or abroad) and save a fraction
$s$ of gross national income (wages plus the gross return on wealth): $\dot W=s\,\mathrm{GNI}-\delta W$. Numbers:
$\alpha=1/3$, $r^*=0.04$, $\delta=0.05$, $n+g=0.03$.}

\textbf{a)} (3 points) Characterize GDP per effective worker and its dynamics. What does a rise in $s$ do to GDP?
What does this imply for the Lucas paradox?
\resolespaco
\begin{sol}
\passo{capital is pinned}{Firms rent capital at $r^*+\delta$: $f'(k)=\alpha k^{\alpha-1}=r^*+\delta$ with $k=K/(AhL)$, so
$k^*=\big(\frac{\alpha}{r^*+\delta}\big)^{1/(1-\alpha)}$ at every date: capital flows in or out instantly.}
\passo{GDP}{$y^*=k^{*\alpha}$ is constant, so GDP per worker $Ahy^*$ grows at $g$ from the start: no transition, and $s$
does not enter. A higher $s$ changes ownership of capital, not production.}
\passo{Lucas paradox}{With equal $A$ and $h$, the closed-economy MPK gap would be several-fold (item 3(a)) and
capital would pour in. Small observed flows mean that $A$ and $h$, which set $k^*$ through $Ah$, differ.}
\refbib{Romer (2012), problem 4.4; Kurlat (2020), \S5.3.}
\rubrica{
\item $k^*$ from the world rate: 1 pt.
\item No transition and no role for $s$ in GDP: 1 pt.
\item Lucas-paradox reading: 1 pt.
}
\end{sol}

\textbf{b)} (3 points) Derive the law of motion of wealth per effective worker, its steady state and the condition
for stability.
\resolespaco
\begin{sol}
\passo{national income}{Wages per effective worker are $(1-\alpha)y^*$; wealth earns $r^*+\delta$ gross. With $\omega\equiv W/(AhL)$:
$\mathrm{GNI}/(AhL)=(1-\alpha)y^*+(r^*+\delta)\omega$.}
\passo{law of motion}{$\dot\omega/\omega=\dot W/W-n-g$ gives
\[\dot\omega=s(1-\alpha)y^*-\Lambda\,\omega,\qquad\Lambda\equiv n+g+\delta-s(r^*+\delta).\]}
\passo{solution}{$\omega(t)=\omega^*+(\omega(0)-\omega^*)e^{-\Lambda t}$, $\omega^*=\frac{s(1-\alpha)y^*}{\Lambda}$: it converges iff $\Lambda>0$. If $\Lambda\le0$
wealth per effective worker grows without bound and the economy eventually ceases to be small.}
\refbib{Romer (2012), \S1.3 (same method).}
\rubrica{
\item GNI per effective worker: 1 pt.
\item Law of motion: 1 pt. Omitting the return on wealth: max.\ 1 pt in the item.
\item Explicit solution and stability: 1 pt.
}
\end{sol}

\textbf{c)} (3 points) Find the saving rate above which Brazil is a net creditor in the steady state and interpret
it. Find GNI/GDP in the steady state in closed form.
\resolespaco
\begin{sol}
\passo{creditor condition}{Net foreign assets are $\omega^*-k^*$. Using $(r^*+\delta)k^*=\alpha y^*$:
$\omega^*>k^*\iff s(1-\alpha)y^*>k^*\Lambda=k^*(n+g+\delta)-s\alpha y^*\iff s\,y^*>(n+g+\delta)k^*$, i.e.
\[s>\bar s\equiv\frac{(n+g+\delta)k^*}{y^*}=\frac{\alpha(n+g+\delta)}{r^*+\delta},\]
the closed-economy saving rate whose steady state is $k^*$.}
\passo{GNI/GDP}{$\frac{\mathrm{GNI}}{\mathrm{GDP}}=(1-\alpha)+\alpha\frac{\omega^*}{k^*}$ and $\frac{\omega^*}{k^*}=\frac{s(1-\alpha)(r^*+\delta)}{\alpha\Lambda}$, so
\[\frac{\mathrm{GNI}}{\mathrm{GDP}}=(1-\alpha)\Big[1+\frac{s(r^*+\delta)}{\Lambda}\Big]=\frac{(1-\alpha)(n+g+\delta)}{\Lambda},\]
which equals 1 exactly at $s=\bar s$: creditor status and GNI above GDP are the same statement.}
\refbib{Romer (2012), \S1.5 and problem 4.4.}
\rubrica{
\item Creditor condition derived: 1.5 pts.
\item GNI/GDP in closed form and the link at $\bar s$: 1.5 pts.
}
\end{sol}

\textbf{d)} (3 points) With $s=0.15$ compute $k^*$, $y^*$, $\omega^*$, net foreign assets, the share of domestic capital owned
abroad, GNI/GDP and the half-life of wealth. Draw $\dot\omega$ against $\omega$ for $s=0.15$ and $s=0.25$, marking $k^*$, and
describe the transition after $s$ rises.
\resolespaco
\begin{sol}
\passo{numbers}{$k^*=(\tfrac{1/3}{0.09})^{3/2}=7.128$, $y^*=1.925$, wages $1.283$; $\Lambda=0.0665$; $\omega^*=0.1925/0.0665=2.894$; net
foreign assets $-4.234$ (59\% of domestic capital foreign-owned); GNI/GDP $=0.0533/0.0665=0.802$; half-life
$\ln2/0.0665=10.4$ years.}
\begin{center}
\begin{tikzpicture}
\begin{axis}[width=10cm,height=5.5cm,xmin=0,xmax=8,ymin=-0.2,ymax=0.36,axis lines=middle,
  xlabel={$\omega$},ylabel={$\dot\omega$},tick label style={font=\footnotesize},label style={font=\small},
  legend style={font=\footnotesize},legend pos=north east]
\addplot[thick,blue,domain=0:8]{0.19245-0.0665*x};
\addplot[thick,red,domain=0:8]{0.32075-0.0575*x};
\addplot[dotted,thick] coordinates {(7.128,-0.2) (7.128,0.36)};
\legend{$s=0.15$,$s=0.25$}
\node[font=\footnotesize] at (axis cs:2.894,-0.03) {$\omega^*_{.15}$};
\node[font=\footnotesize] at (axis cs:5.578,-0.03) {$\omega^*_{.25}$};
\node[font=\footnotesize,anchor=west] at (axis cs:7.15,0.3) {$k^*$};
\end{axis}
\end{tikzpicture}
\end{center}
\passo{transition}{Each line crosses zero at its steady state and slopes down (slope $-\Lambda$): stable. After the rise
in $s$, $\omega$ climbs from $2.894$ toward $5.578$ (half-life $\ln2/0.0575=12.1$ years), still left of $k^*$: GNI rises
gradually while GDP does not move.}
\refbib{Romer (2012), problem 4.4.}
\rubrica{
\item Numbers: 1.5 pts.
\item Graph: axes, two downward lines with the right intercepts and zeros, $k^*$ marked: 1 pt.
\item Transition read off the graph: 0.5 pt.
}
\end{sol}

\textbf{e)} (3 points) \emph{Parameter analysis.} How do $\omega^*/k^*$, GNI/GDP and the half-life of wealth respond to $s$, $r^*$,
$n+g$ and $\alpha$? Prove the signs, fill the table for the baseline and for $s=0.25$, $r^*=0.06$, $n+g=0.02$, $\alpha=0.4$
(one change at a time), and give the limits as $s\to0$ and as $s$ approaches its upper bound.
\resolespaco
\begin{sol}
\passo{signs}{From (c), $\frac{\omega^*}{k^*}=\frac{s(1-\alpha)(r^*+\delta)}{\alpha\Lambda}$, $\frac{\mathrm{GNI}}{\mathrm{GDP}}=\frac{(1-\alpha)N}{\Lambda}$, half-life $=\frac{\ln2}{\Lambda}$, with $N=n+g+\delta$ and
$\Lambda=N-s(r^*+\delta)$. $\Lambda$ falls with $s$ and $r^*$ and rises with $N$; it does not depend on $\alpha$. Hence: all three rise
with $s$ and with $r^*$ (a higher return compounds wealth faster); $\omega^*/k^*$ and GNI/GDP fall with $n+g$
($\partial[N/(N-a)]/\partial N=-a/(N-a)^2<0$) and so does the half-life; $\omega^*/k^*$ and GNI/GDP fall with $\alpha$ (labor income,
the source of saving, is a smaller share), while the half-life does not depend on $\alpha$.}
\passo{table}{
\begin{center}
\begin{tabular}{lccc}
\toprule
 & $\omega^*/k^*$ & GNI/GDP & half-life (years) \\
\midrule
baseline & 0.406 & 0.802 & 10.4 \\
$s=0.25$ & 0.783 & 0.928 & 12.1 \\
$r^*=0.06$ & 0.520 & 0.840 & 10.9 \\
$n+g=0.02$ & 0.478 & 0.826 & 12.3 \\
$\alpha=0.4$ & 0.305 & 0.722 & 10.4 \\
\bottomrule
\end{tabular}
\end{center}}
\passo{limits}{$s\to0$: $\omega^*\to0$ and GNI/GDP $\to1-\alpha$ (residents own nothing and earn only wages). $s\to N/(r^*+\delta)$
($=0.89$ here): $\Lambda\to0$, wealth, GNI/GDP and the half-life explode; the small-economy assumption fails before.}
\intuicao{Openness separates where capital works (set by productivity and the world rate) from who owns
it (set by saving); every parameter that raises the return on owned wealth relative to dilution by growth
makes residents richer without changing what Brazil produces.}
\refbib{Romer (2012), problem 4.4; Kurlat (2020), \S1.1 (GDP vs GNI).}
\rubrica{
\item Signs proved from the closed forms: 1.5 pts.
\item Table: 1 pt.
\item Limits: 0.5 pt.
}
\end{sol}

\end{document}
